\documentclass[11pt]{article}
\usepackage[margin=0.85in]{geometry}
\usepackage[T1]{fontenc}
\usepackage[utf8]{inputenc}

\setlength{\parindent}{0pt}
\setlength{\parskip}{0.55em}
\makeatletter
\providecommand{\@mlabel}[2]{}
\makeatother

\begin{document}

\begin{flushright}
Centre for Computational Science and Mathematical Modelling\\
Coventry University\\
Coventry, United Kingdom\\
olarinoyem@coventry.ac.uk\\
\today
\end{flushright}

The Editors\\
\textit{Artificial Intelligence in Medicine}\\
Elsevier

Dear Editors,

On behalf of my co-authors, I am pleased to submit our manuscript, ``MSAGAT-Net: Multi-Scale Temporal Adaptive Graph Attention for Efficient Spatiotemporal Epidemic Forecasting,'' for consideration as an Original Research Article in \textit{Artificial Intelligence in Medicine}.

The manuscript introduces MSAGAT-Net, a multi-scale adaptive graph attention network for spatiotemporal epidemic forecasting. The method addresses important limitations of existing graph neural forecasting models, including reliance on predefined adjacency matrices, fixed spatial aggregation depths, and reduced stability at longer forecast horizons. Its main methodological contributions are a graph attention mechanism that combines content-based scores with a learnable low-rank graph bias and an optional adjacency prior, an adaptive multi-hop spatial diffusion module whose depth scales with graph size, and a progressive prediction refinement module for multi-horizon forecasting.

We believe the manuscript is well aligned with the scope of \textit{Artificial Intelligence in Medicine} because it combines methodological advances in AI with a clear healthcare application. We evaluate MSAGAT-Net on real-world forecasting tasks spanning influenza, COVID-19, and ICU bed occupancy, demonstrating relevance to epidemic preparedness, population health surveillance, and healthcare resource planning. The model achieves strong performance against established baselines, including RMSE reductions of up to 23.5\% on LTLA-COVID forecasting and 22.2\% on NHS-ICUBeds forecasting. Ablation analyses across three graph scales further clarify how each component contributes across forecast horizons and graph sizes.

This manuscript is original, has not been published elsewhere, and is not under consideration by another journal. All authors have approved the manuscript and agree with its submission to \textit{Artificial Intelligence in Medicine}. The authors declare no competing financial interests or personal relationships that could have appeared to influence the work reported in this paper.

Thank you for your time and consideration.

\begin{flushright}
Yours sincerely,

\vspace{0.9em}

Michael Ajao-olarinoye\\
Centre for Computational Science and Mathematical Modelling\\
Coventry University, Coventry, United Kingdom\\
olarinoyem@coventry.ac.uk
\end{flushright}

\end{document}
